\section{Introduction}

    \subsection{Background}
        \begin{frame}{Background}
            \begin{itemize}
                \item Increased interest in civilian applications.
                \item For many applications, a small scale vehicle is desired.
                \item MAV - Micro Air Vehicle; UAV weighing $5$~kg or less\footnote{Definitions differ}.
            \end{itemize}
            \centering
            \begin{minipage}{\linewidth}
                \centering
                \sfig{0.3}{skeldar}
                \sfig{0.3}{microUAV}
                \sfig{0.3}{rmax}
            \end{minipage}
        \end{frame}
        \note{
            First a bit of background;

            During the past few years, there has been a significant increase
            of interest in the field of unmanned aerial vehicles.
            While the military industry is the major developer of UAVs,
            the maturity and availability of technology has enabled applications
            in for example search-and-rescue operations and even hobby-level implementations.

            Many UAVs are of considerable size to carry load and to increase flight-time.
            As new applications develop however, the interest has increased in small scale UAVs
            such as the LinkQuad, towards which this thesis have been targeted.
            This type of UAVs are known as MAVs, Micro Air Vehicles.
        }

    \subsection{The LinkQuad Platform}
        \begin{frame}{The LinkQuad Platform}
            Quadrotor research platform developed at AIICS at the
            Department of Computer and Information Science of Linköping University.

            Available sensors: Accelerometers, gyroscopes, pressure sensor, camera, magnetometers, GPS...

            \fig{0.5}{linkquadoutside}
        \end{frame}
        \note{
            The LinkQuad platform, which you can see here, is a quadrotor developed here at IDA.
            It is equipped with a standard set of inertial sensors - accelerometers and gyroscopes -
            as well as pressure sensors, magnetometers and GPS.
            It can also be optionally equipped with a camera, which was used in this thesis.
        }
        \begin{frame}{The LinkQuad Platform}
            \begin{itemize}
                \item Dual gumstix micro-computers (Linux).
                \item Sensor-board with dual microcontrollers for sensor sampling, data logging and low-level control.
            \end{itemize}
            \begin{tikzpicture}[auto,>=latex']
            \tikzstyle{mcu} = [rectangle, draw=none, rounded corners=1mm, fill=Emerald,
                text centered, anchor=north, text=white,font=\small]
            \tikzstyle{sensor} = [rectangle, draw=none, fill=black,
                text centered, anchor=north, text=white,font=\tiny, node distance=0.8cm]
            \tikzstyle{camera} = [rectangle, draw=none, rounded corners=1mm, fill=red,
                text centered, anchor=north, text=white,font=\small]
            \tikzstyle{framework} = [rectangle, draw, fill=OliveGreen,
                text centered, anchor=north, text=white,font=\small,thin, minimum width=2cm,node distance=0.7cm]
            \tikzstyle{gumstix} = [rectangle, draw=none, fill=Green, rounded corners=2mm,
                text centered, anchor=north, text=white,font=\small]
            \tikzstyle{gumstixtext} = [rectangle, draw=none,
                text centered, align=center, anchor=north, text=white,font=\small]

            \node [gumstixtext] (primarytext) {Primary computer \\ \tiny{State estimation and control}};
            \node [framework, below of=primarytext] (crap) {CRAP};
            \node [gumstixtext, right of=primarytext, node distance=5cm] (secondarytext) {Secondary computer \\ \tiny{Video processing}};
            \node [framework, below of=secondarytext] (ptam) {PTAM};

            \begin{pgfonlayer}{background}
                \node [gumstix, fit=(primarytext) (crap)] (primary) {};
            \end{pgfonlayer}
            \begin{pgfonlayer}{background}
                \node [gumstix, fit=(secondarytext) (ptam)] (secondary) {};
            \end{pgfonlayer}

            \node [mcu, below of=primary, xshift=-1.5cm, node distance=1.5cm] (smcu) {Sensor MCU};
            \node [mcu, right of=smcu, node distance=3cm] (cmcu) {Control MCU};
            \node [sensor, below of=smcu, xshift=-1.2cm] (accelerometer) {Acc.};
            \node [sensor, right of=accelerometer] (gyro) {Gyro};
            \node [sensor, right of=gyro] (magnetometer) {Mag.};
            \node [sensor, right of=magnetometer] (pressure) {hPa};
            \node [camera, below of=secondary, node distance=1.5cm] (camera) {Camera};
            \node [sensor, below of=cmcu] (esc) {Speed controller};
            \node [sensor, right of=esc, node distance=1.8cm] (sdcard) {SD-card};

            \path[-]
                        (primary) edge (secondary)
                        (primary.south) edge (smcu.north)
                        (primary.south) edge (cmcu.north)
                        (cmcu.south) edge (esc.north)
                        (smcu.south) edge (accelerometer.north)
                        (smcu.south) edge (gyro.north)
                        (smcu.south) edge (magnetometer.north)
                        (smcu.south) edge (pressure.north)
                        (secondary) edge (camera)
                        (smcu.east) edge (cmcu.west)
                        (cmcu.south) edge (sdcard.north)
                        ;
            \end{tikzpicture}
        \end{frame}
        \note{
            One particularly notable feature on the LinkQuad is the dual
            gumstix Linux microcomputers. The fact that there are two allows us to dedicate one
            each of these to video processing and state estimation respectively.

            These gumstix computers measure only centimeters in size, yet are
            quite powerful in terms of computational power. They do, however,
            lack floating point processor units, which in the end was a
            major limiting factor as to why only simulational flight was achieved.

            Aside from the gumstix microcomputers, the LinkQuad is also
            equipped with microcontrollers for sampling the sensors, logging and
            actuating control signals.
        }

    \subsection{Problem Formulation}
        % Why was a high-level filtering and control framework needed?
        \begin{frame}[c]{Problem Formulation}

            \textbf{Primary goal:} Develop a control system
                for the LinkQuad that use video-based positioning
                to enable stable landing.

            \textbf{Breakdown:}
            \begin{itemize}
                \item Localization with Video-based SLAM.
                \item Sensor Fusion with available sensors.
                \item Use state estimate for control.
                \item Generate control reference for landing procedure.
                \item Landing detection.
            \end{itemize}
        \end{frame}
        \note{
            To be able to act independently, a UAV should be able to land autonomously.
            Therefore, the goal of this thesis was
            to develop a control system that, with the help of camera positioning,
            could be used to land the LinkQuad.

            To accomodate and use the quite complex measurements
            reeived from the camera, the approach has been to develop
            and implement high-level algorithms for state estimation and control,
            as well as to extend previous work on video-based positioning to
            autonomous operation.
            A reference trajectory was then generated to perform the landing
            procedure until a landing was detected.
        }

    \subsection{Method}
        \begin{frame}{Method: Video-based Positioning}
            Extract 3D information from 2D video-stream.
            \begin{itemize}
                \item vSLAM - Visual SLAM.
                \item Simultaneous Localisation And Mapping to track features in the video-stream.
                \item Gives orientation and position relative to the features.
                \item High computational demands.
                \item Complex sensor measurements.
                %~ \item Fits nicely into a state-space filtering framework.
            \end{itemize}

            \begin{figure}[H]
            \centering
                \includegraphics[trim = 0mm 30mm 30mm 30mm,clip,width=0.45\textwidth]{\currentchapter/figures/ptamtracking}
                \hspace{0.05\textwidth}
                \includegraphics[trim = 20mm 30mm 10mm 30mm,clip,width=0.45\textwidth]{\currentchapter/figures/ptamtracking}
            \end{figure}
        \end{frame}
        \note{
            One of the main technologies that were studied in this thesis
            is video based SLAM. That is, Simultaneous Localization - of the vehicle - And Mapping - of the environment - using
            video-stream data.

            There are several vSLAM algorithms which
            work in similar ways - by extracting trackable features from the video images,
            and tracking their changes in position between the video frames.
            These changes are used to estimate the camera's position relative to the
            features.

            Video processing generally has high computational demands, and
            the resulting measurements are quite complex to interpret.
            It was nescessary to implement a high-level filtering framework to
            fuse the video positioning with the state estimation of the quadrotor.
        }

        \begin{frame}{Method: Sensor Fusion}
            Current implementation is based on
            complementary filtering.
            \scriptsize
            \begin{equation*}
                \text{angle} = \mathcal{H}_{HP}\left( \int \text{gyro} \, \text{dt} \right) + \mathcal{H}_{LP}\left( accelerometers \right)
            \end{equation*}
            \normalsize
            Performance is adequate, but the c.f is difficult to extend.

            Also, camera measurements
            \begin{itemize}
                \item Cannot be used directly in the current filter,
                \item Fits well into a standard state-space filter framework.
            \end{itemize}

            A high-level filtering framework with an advanced motion model
            was implemented and applied for state estimation.
        \end{frame}
        \note{
            The current observer - that is, implementation of state estimation - is based on the
            complementary filter, and while the performance of the filter
            is quite adequate for its current purpose, adding more sensors - especially
            such as the camera positioning - is quite difficult.

            The de facto standard techique for nonlinear state estimation,
            the Extended Kalman Filter, was ultimately applied for the state estimation.
            Extensive work was in fact put into creating a physical motion model,
            and while a simpler model probably would have sufficed for the state estimation,
            the same model could thereby be re-used later, for instance in the nonlinear
            control that is proposed in the thesis.
        }

        \begin{frame}{Method: Control}
            With a linear motion model available, control signals can be computed optimally.
            Linear Quadratic control offers:
            \begin{itemize}
                \item Simple implementation.
                \item Simple tuning.
                \item Optimal control.
            \end{itemize}
            Motion model needs to be linear, which is not the case for a quadrotor.
            This constraint can be circumvented by an extension to LQ control using the
            \textbf{State-Dependent Riccati Equation}.
        \end{frame}
        \note{
            With a motion model available, the Linear Quadratic controller was near
            at hand to be applied. The problem is, the LQ requires the motion model to be linear.

            In the thesis, I study a technique to bypass this constraint
            called State-Dependent Riccati Equations, which basically solves a modified
            problem repeatedly online, extending the LQ methodology
            using a trick I will present later.

            This proposed trick makes it possible to regain much of the simplicity of
            the linear case, while extending the utility to a far more general case.
        }

    \tikzstyle{state} = [rectangle, draw=none, rounded corners=1mm, fill=OliveGreen,
                    text centered, anchor=north, text=white, minimum width=2cm, minimum height=1cm, node distance=2cm, align=center]
    \tikzstyle{condition} = [diamond, draw=none, fill=red!80,
                    text centered, anchor=north, text=white, minimum height=1cm, node distance=2.3cm, align=center]
    \tikzstyle{choice}=[font=\scriptsize]

        \begin{frame}{Method: Reference Generation}
            With properly implemented control, landing is a matter of descending
            steadily until landing is detected.
            \begin{figure}[h]
            \footnotesize\centering
                \begin{tikzpicture}[auto]
                    \node (entry) {Entry};
                    \node [state, right of=entry] (halt) {Halt to hover\\ {\tiny Set $V_{ref}=0$}};
                    \node [condition,right of=halt] (halted) {$|V| < \epsilon$};
                    \node [state, below of=halt, node distance=1.8cm] (descend) {Descend \\ {\tiny Set $V_{ref,z} > 0$}};
                    \node [condition,right of=descend] (landed) {Landed};
                    \node [state, below of=descend, node distance=1.8cm] (spindown) {Spin down};
                    \path[->]
                        (entry.east) edge (halt.west)
                        (halt.east) edge (halted.west)
                        (halted) edge[loop right,in=90,out=-135] (descend.north)
                        (descend.east) edge (landed.west)
                        (landed) edge[loop right,in=90,out=-135] (spindown.north)
                    ;
                \end{tikzpicture}
            \end{figure}
        \end{frame}
        \note{
            I also implemented a state-machine engine to generate
            reference signals for the controlled descent and landing of the quadrotor.
            The basic idea is to use the positioning for a stabilized descent
            and to detect the event of landing when the vehicle has hit the ground.

            On the screen, you can see the basic outline of the landing scheme;
            First, the quadrotor is put into a halt, hovering in mid-air.
            The vehicle then starts its descent which will continue until a
            landing has been detected.
        }

        \begin{frame}{Method: Landing Detection}
            Landing is generally associated with a lack of descent.

            \vspace{0.7cm}

            \textbf{Two interesting states to monitor:}
            \begin{itemize}
                \item Vertical velocity.
                \item Vertical wind velocity.
            \end{itemize}

            \vspace{0.7cm}
            Without modeled ground force, the observer will explain the
            added force with upward winds.

            %~ Beause only simulated landing was performed in this thesis,
            %~ the implementation needs further work.
        \end{frame}
        \note{
            When landing occurs, the observer's current motion model will try
            to describe this sudden loss of vertical velocity and will,
            after a while's settling time, put the blame on
            an increased upward force from the modeled wind.

            To detect landing, it is thus of interest to monitor both
            the vertical velocity and the estimated wind velocity
            for increased detection reliability.
        }

%~ \begin{frame}{The LinkQuad Platform}
    %~ \fig{LinkQuad}{0.3}
%~ \end{frame}
